\section{A limitation of pairing concentration}
\label{sec:random-seven-concentration-counterexample}

Vanishing terminal-path densities do not force four-terminal forest mass into one pairing.
The following example explains why a restriction on cyclic cuts is necessary.

\begin{proposition}
\label{prop:random-seven-concentration-counterexample}
For every integer $b\ge2$, there is a connected simple normalized four-pole $X_b$ of rank $h=4b+1$ such that
\[
 F_{\{1,2,3,4\}}(X_b)=2^h,
 \qquad
 f_{12|34}(X_b)=f_{14|23}(X_b)=2^{h-1},
 \qquad
 f_{13|24}(X_b)=0,
\]
and
\[
 P_{ij}(X_b)=2^{2b+1}
 \qquad(1\le i<j\le4).
\]
In particular, all normalized path densities tend to zero, while two forest-pairing densities remain equal to $1/2$.
\end{proposition}

\begin{proof}
Let $Q_b$ be a saturated simple two-pole of rank $b$.
For example, use the capped singleton ladder with $b-1$ rung positions.
It has $2^b$ simple paths between its two terminals.

Start with a cycle $v_1v_2v_3v_4v_1$.
For each $i$, take a disjoint copy of $Q_b$, with terminals $a_i,t_i$.
Add the bridge $v_ia_i$.
Use $t_1,t_2,t_3,t_4$ as the four ports.
Every other vertex now has degree three.
Bridges do not change cycle rank, so the resulting rank is $h=1+4b$.

A join with odd boundary at all four ports must use all four bridges.
Its restriction to each $Q_b$ is an $a_i$--$t_i$ path, since that two-pole is saturated.
The restriction to the central four-cycle has odd degree at each vertex.
There are exactly two possibilities: its two alternating matchings.
Each gives a forest and has $2^{4b}$ independent choices in the four attached two-poles.
The stated forest counts follow.

A simple $t_i$--$t_j$ path uses exactly the two attached poles at $i,j$.
It cannot enter another attached pole, because that would require traversing its bridge twice.
There are $2^b$ choices in each used pole and two arcs between $v_i,v_j$ in the central cycle.
Thus its count is $2^{2b+1}$, as required.
\end{proof}

Consequently, no absolute constant $C$ can satisfy
\[
 f_\alpha(X)f_\beta(X)
 \le C\,2^{\beta(X)}\sum_{i<j}P_{ij}(X)
 \qquad(\alpha\ne\beta)
\]
for all normalized four-poles.
For the displayed family, the ratio of the two sides without $C$ is $2^{2b}/24$.

Each attached pole gives a cyclic two-edge cut in every cubic completion.
Its boundary consists of its attachment bridge and its filled port edge.
Therefore this example leaves the corresponding concentration question open for graphs without small cyclic cuts.
